\renewcommand{\headrulewidth}{0pt}

\begin{center}
{\Large\bfseries ХУРААНГУЙ}
\end{center}
\addcontentsline{toc}{chapter}{Хураангуй}

\vspace{0.5cm}

\begin{center}
{\fontsize{14pt}{16pt}\selectfont \textbf{Камерын 2D дүрсээс машин сургалтаар 3D загвар үүсгэх аргын судалгаа}}
\end{center}

\begin{flushright}
\itshape
\textbf{Оюутан:} Б.Эрдэнэбаяр\\
\textbf{Удирдагч багш:} Н.Соронзонболд
\end{flushright}

\vspace{0.5cm}

\noindent\textbf{Судалгааны үндэслэл:}
3D загварчлал нь компьютер график, роботикийн ертөнц, VR/AR болон соёлын өвийн дижитал хадгалалт зэрэг олон салбарт чухал. Уламжлалт Structure-from-Motion (SfM) болон Multi-View Stereo (MVS) аргууд нь геометрийн үнэн зөв байдал, тогтвортой байдлаараа давуу ч, тооцоолол өндөр шаардлага ба гэрэлтэлт, зурагнуудын давхцал бага байдалд сул талтай. Энэ судалгаа нь COLMAP-д суурилсан практик pipeline-ийг сайжруулж, машин сургалтын аргаар зарим үе шат (feature extraction, outlier filtering, confidence scoring) илүү үр дүнтэй болгох боломжийг судална.

\vspace{0.3cm}

\noindent\textbf{Судалгааны зорилго:}
Олон өнцгөөс авсан 2D зурагнаас практик, тогтвортой 3D модель (геометр + текстур) үүсгэх pipeline боловсруулахын тулд:
\begin{itemize}
\item Зургийн урьдчилсан боловсруулалт, камерын калибраци, feature extraction-г автоматжуулах
\item COLMAP ашиглан эхний sparse ба dense point cloud үүсгэх
\item Үр дүнг геометрийн үнэн зөв (reprojection error, Chamfer distance), бүрэн байдал (completeness) ба визуал чанараар үнэлэх
\item Уламжлалт MVS болон SfM-ын оронд Машин сургалтын аргуудыг орлуулах боломжийг практик аргаар судлах
\end{itemize}

\vspace{0.3cm}

\noindent\textbf{Арга зүй:}
Pipeline-д дараах үе шатууд орно: зурагнуудын kinematic болон photometric урьдчилсан боловсруулалт, SIFT болон SuperPoint-аас гарсан онцлогуудыг харьцуулах, COLMAP-н SfM ба MVS модулиудыг ажиллуулах. ML-аргыг (жиш: small CNN эсвэл Random Forest) feature descriptor-уудын чанарыг үнэлж, таарууарчихааг (mismatches) шүүн ялгах болон таамаглалын итгэлцлийг тодорхойлоход ашиглав. Моделийн гүйцэтгэлийг synthetic болон real-world dataset-ууд дээр харьцуулсан.

\vspace{0.3cm}

\noindent\textbf{Үр дүн:}
Анхны туршилтуудын дагуу COLMAP-н уламжлалт pipeline нь геометрийн нарийвчлал сайн хадгалсан. ML-н тусламжтай feature filtering ба confidence scoring оруулснаар outlier-ууд буурч, reprojection error-ыг бага хэмжээгээр сайжруулсан бөгөөд reconstruction-ийн тогтвортой байдал гэрэлтэлтийн хувирал үүссэн тохиолдолд сайжирсан. Мөн reconstructed mesh-ийн бүрэн байдал болон texture continuity дээр давуу тал ажиглагдсан.

\vspace{0.3cm}

\noindent\textbf{Дүгнэлт:}
SfM + MVS (COLMAP)-д тулгуурласан pipeline нь практик 3D reconstruction-д хамгийн боломжит суурь хэвээр байна. Машин сургалтын аргуудыг хэсэгчилсэн байдлаар (feature quality, outlier removal, confidence scoring) нэвтрүүлэх нь reconstruction-ийг тогтвортой, практик нөхцөлд илүү найдвартай болгодог. Үргэлжлүүлэх ажилд hybrid сайжруулалтуудыг олон төрлийн зураглалтай, real-time ойролцоо шийдлүүд рүү шилжүүлэх боломжтой.

\vspace{0.3cm}

\noindent\textbf{Түлхүүр үг:}
3D Reconstruction, COLMAP, Structure-from-Motion, Multi-View Stereo, Feature Matching, Machine Learning, Computer Vision.

\clearpage

% ========================= ABSTRACT ========================= %

\begin{center}
{\Large\bfseries ABSTRACT}
\end{center}
\addcontentsline{toc}{chapter}{Abstract}

\vspace{0.5cm}

\begin{center}
{\fontsize{14pt}{16pt}\selectfont \textbf{3D Reconstruction from 2D Camera Images Using Machine Learning}}
\end{center}

\begin{flushright}
\itshape
\textbf{Student:} B. Erdenebayar\\
\textbf{Supervisor:} N. Soronzonbold
\end{flushright}

\vspace{0.5cm}

\noindent\textbf{Background:}
3D reconstruction is essential for graphics, robotics, VR/AR, and digital preservation. Classical Structure-from-Motion (SfM) and Multi-View Stereo (MVS) methods provide reliable geometry but are computationally demanding and affected by lighting changes and low image overlap. This thesis focuses on improving a COLMAP-based SfM/MVS pipeline with targeted machine-learning enhancements to increase robustness and practical usability.

\vspace{0.3cm}

\noindent\textbf{Objective:}
Design a practical pipeline to reconstruct geometry and texture from multi-view images by:
\begin{itemize}
\item Automating preprocessing, camera calibration, and feature extraction
\item Producing initial sparse and dense point clouds using COLMAP
\item Introducing ML-based feature filtering, outlier detection, and confidence scoring
\item Evaluating reconstruction with reprojection error, Chamfer distance, completeness, and visual quality
\end{itemize}

\vspace{0.3cm}

\noindent\textbf{Methodology:}
The pipeline includes photometric and geometric preprocessing, feature detection (SIFT, SuperPoint), matching, and COLMAP SfM/MVS processing. Lightweight ML models (e.g., small CNNs or tree-based classifiers) are trained to assess descriptor quality and filter mismatches. Experiments run on both synthetic and real-world datasets for quantitative and qualitative comparison.

\vspace{0.3cm}

\noindent\textbf{Results:}
Baseline COLMAP reconstructions show high geometric accuracy. Integrating ML-based filtering and confidence scoring reduced outliers and improved reconstruction stability under varying lighting and limited overlap conditions. Improvements were observed in reprojection error, completeness, and texture continuity.

\vspace{0.3cm}

\noindent\textbf{Conclusion:}
A COLMAP-based SfM/MVS pipeline augmented with targeted machine-learning components yields more robust and practical 3D reconstructions for real-world scenarios. Future work will explore further hybridization, runtime optimization, and extensions toward near real-time reconstruction.

\vspace{0.3cm}

\noindent\textbf{Keywords:}
3D Reconstruction, COLMAP, Structure-from-Motion, Multi-View Stereo, Feature Matching, Machine Learning, Computer Vision.
